%% kapa: exam=1 round=34 year=2023 period=1 form=A subject=real_estate subject_ko=부동산학원론 seq=35 number=115
%% source: Q-Net 감정평가사 기출문제 (cst00309) artlSeq=5211789 file=raw/1008/5211789/2023년도 제34회 감정평가사 1차시험 1교시.pdf
%% conversion: kordoc 4.24.51 -> tools/gen1.py (auto); review=proofread-sonnet; flags=none
\begin{question}{115}
감정평가에 관한 규칙과 감정평가 실무기준 상 임대료 감정평가에 관한 설명으로 옳지 않은 것은?
\begin{choices}
  \item 임대사례비교법으로 감정평가 할 때 임대사례에 특수한 사정이나 개별적 동기가 반영되어 수집된 임대사례의 임대료가 적절하지 못한 경우에는 사정보정을 통해 그러한 사정이 없었을 경우의 적절한 임대료 수준으로 정상화하여야 한다.
  \item 시점수정은 대상물건의 임대료 변동률로 함을 원칙으로 한다.
  \item 감정평가법인등은 임대료를 감정평가할 때에 임대사례비교법을 적용해야 한다.
  \item 적산법은 원가방식에 기초하여 대상물건의 임대료를 산정하는 감정평가방법이다.
  \item 수익분석법이란 일반기업 경영에 의하여 산출된 총수익을 분석하여 대상물건이 일정한 기간에 산출할 것으로 기대되는 순수익에 대상물건을 계속하여 임대하는 데에 필요한 경비를 더하여 대상물건의 임대료를 산정하는 감정평가방법을 말한다.
\end{choices}
\end{question}
\answer{2}
